\documentclass[10pt,a4paper]{article}
\usepackage[margin=0.65in]{geometry}
\usepackage[T1]{fontenc}
\usepackage[utf8]{inputenc}
\usepackage[english]{babel}
\usepackage[expansion=false,protrusion=true]{microtype}
\usepackage{newtxtext}
\usepackage{newtxmath}
\usepackage{graphicx}
\usepackage{booktabs}
\usepackage[table]{xcolor}
\usepackage{tabularx}
\usepackage{array}
\usepackage{float}
\usepackage[ruled,vlined]{algorithm2e}
\usepackage[hidelinks]{hyperref}

\graphicspath{{docs_second/figures/}}

\pagenumbering{gobble}
\setlength{\parindent}{0pt}
\setlength{\parskip}{5pt}

\definecolor{iacBlue}{RGB}{31,78,121}
\definecolor{iacGray}{RGB}{245,245,245}

\newcommand{\blocktitle}[1]{\textbf{\textcolor{iacBlue}{#1}}}
\setlength{\fboxsep}{6pt}
\setlength{\fboxrule}{0.6pt}

\begin{document}

{\LARGE \textbf{ARTPS: Operationally Safe Target Prioritization via Depth-Enhanced Hybrid Anomaly Detection}}\par
\vspace{2pt}
{\normalsize \textbf{Poyraz BAYDEMİR} \;\;|\;\; \textit{Selçuk University (Konya, Türkiye)}}\par
{\small \textbf{Category:} Science \& Exploration / Space Exploration Symposium}\par
\vspace{6pt}
\hrule
\vspace{8pt}

\begin{minipage}[t]{0.60\linewidth}
\blocktitle{Problem (Why).}\;
Planetary rovers operate under strict bandwidth and energy budgets with delayed ground feedback, forcing more autonomy onboard. In field imagery, visual-only target selection is brittle: illumination changes, cast shadows, and low-texture regions can create false positives, while visually similar rock targets can suppress high-value candidates (false negatives) under naive diversity policies.

\blocktitle{Solution (What --- ARTPS).}\;
We introduce ARTPS, an edge-oriented framework that couples monocular depth estimation with multi-cue anomaly detection and robust localization. ARTPS fuses image cues (texture/edges and photometric indicators), depth cues (discontinuities and proximity weighting), and optional learned anomaly maps via an entropy-weighted dynamic fusion mechanism that adapts to field conditions.

\blocktitle{Novelty (How it becomes operationally safe).}\;
Beyond accuracy, ARTPS targets reliability in decision-making. We combine feature-space clustering (latent embeddings) for multi-shape management with a soft similarity penalty, and introduce a \emph{Priority Buffer} that stores high-raw-score targets that were down-weighted by diversity policy for second-chance re-evaluation when resources allow.

\blocktitle{Results (Proof).}\;
On NASA Mars rover imagery, ARTPS reaches \textbf{0.894 AUROC}, \textbf{0.847 AUPRC}, and \textbf{0.823 F1}, outperforming a PaDiM/PatchCore baseline (AUROC 0.856). In an edge-oriented configuration (OpenCV preprocessing + fusion + localization, excluding learned depth/AE inference), the core pipeline sustains \textbf{28.1 FPS} at 256\(\times\)256 (12.8 FPS at 384\(\times\)384; 4.0 FPS at 768\(\times\)768) on a development workstation, enabling real-time onboard screening with profile-aware degradation.

\blocktitle{Keywords.}\;
Mars rover; Autonomous Exploration; Anomaly Detection; Computer Vision; Planetary Surfaces.
\end{minipage}\hfill
\begin{minipage}[t]{0.38\linewidth}
\vspace{-2pt}
\fcolorbox{iacBlue}{white}{%
  \begin{minipage}{0.94\linewidth}
  \vspace{2pt}
  {\footnotesize \textbf{Algorithm 1: ARTPS workflow (graphical pseudocode).}}\par
  \vspace{4pt}
  {\scriptsize\ttfamily
  Input: RGB frame I\par
  Output: ranked targets T and buffer B\par
  \vspace{2pt}
  1) D  <- DepthNet(I)\par
  2) I' <- DepthGuidedEnhance(I, D)\par
  3) A\_img, A\_depth, A\_model <- ExtractCues(I', D)\par
  4) A <- EntropyWeightedFusion(A\_img, A\_depth, A\_model)\par
  5) R <- Localize(A)\par
  6) T <- ClusterSelect(R)\par
  7) B <- PriorityBuffer(R)\par
  }
  \vspace{2pt}
  \end{minipage}%
}

\vspace{6pt}
\blocktitle{Key Figures (summary).}\par
\vspace{3pt}
\rowcolors{2}{iacGray}{white}
\begin{tabularx}{\linewidth}{>{\raggedright\arraybackslash}X r}
\toprule
\textbf{Metric} & \textbf{Value} \\
\midrule
AUROC (ARTPS) & 0.894 \\
AUPRC (ARTPS) & 0.847 \\
F1-Score (ARTPS) & 0.823 \\
AUROC (PaDiM/PatchCore) & 0.856 \\
Speed @256\(\times\)256 (core) & 28.1 FPS \\
\bottomrule
\end{tabularx}
\end{minipage}

\vspace{10pt}
\hrule
\vspace{8pt}
\textbf{Impact Statement.}\;
\textbf{This study demonstrates that hybrid AI architectures can enable the next generation of autonomous rovers to perform ``science-on-the-edge'', maximizing scientific return per downlink bit while preserving operational safety.}

\end{document}

